\section{Auxiliary results}\label{app:auxiliary-results}

This appendix collects proofs of the separability, uniqueness, sign, independence, and
choice-implication results stated in the main text.  None is used to prove
Theorem~\ref{thm:main}; separating them keeps Appendices A and B focused on
the representation argument.

\subsection{Separability, uniqueness, and sign}

For $q\in\D(A)$ and $K:A\to\D(O\times R)$, write $\pi^{qK}_{or}(a):=q(a)K(o,r\mid a)/p_{qK}(o,r)$ when
$p_{qK}(o,r)>0$ for the posterior induced by the visible pair, and
$\mu_{qK}:=\sum_{(o,r):p_{qK}(o,r)>0}p_{qK}(o,r)\,\delta_{\pi^{qK}_{or}}$ for its distribution.

\begin{corollary}[No payoff--trace complementarity]\label{cor:matching}
Suppose Axioms \textup{(A1)--(A8)} hold.  Fix $A$ and $q\in\D(A)$.  If channels $K$ and $L$ satisfy
$\E_{qK}[u(O)]=\E_{q,L}[u(O)]$ and $\mu_{qK}=\mu_{qL}$, then $(q,K)\sim(q,L)$.
\end{corollary}

\begin{proof}
The result follows from Theorem~\ref{thm:main} and
$\I_{qK}(A;O,R)=\Sh(q)-\int\Sh(\pi)d\mu_{qK}(\pi)$.
\end{proof}

\begin{proof}[Proof of Proposition~\ref{prop:uniqueness}]
Choose $o_*\in O$ with $u(o_*)=\underline u$.  First note that the range of $V$ is
$[\underline u,\infty)$.  The lower bound is immediate.  For $q\in\D(A)$ and $t\in[0,1]$, let a constant-$o_*$
record channel reveal the realised action with
probability $t$ and otherwise report $*$:
\[
 E_t(o_*,a'\mid a)=t\1_{a'=a},
 \qquad
 E_t(o_*,*\mid a)=1-t.
\]
Then $V(q,E_t)=\underline u+\lambda t\Sh(q)$.  At the uniform lottery on $n$ actions these values fill
$[\underline u,\underline u+\lambda\log n]$; the intervals are unbounded.

For part \textup{(i)}, let $W$ be a common-scale representation and define
$\varphi(v):=W(q,K)$ for any pair with $V(q,K)=v$.  This is well defined: equality of the two $V$-values gives
indifference in their two-block environment by Theorem~\ref{thm:main}, and the common-scale property gives equality
of the two $W$-values.  A strict inequality between $V$-values similarly gives a strict inequality between
$W$-values.  Thus $\varphi$ is strictly increasing and $W=\varphi\circ V$.  The converse is immediate from
Theorem~\ref{thm:main}.

For part \textup{(ii)}, the expected-utility and mutual-information chain rules give, for every
$P:A\to\D(\{1,\ldots,k\})$ and every
$K_1,\ldots,K_k,L_1,\ldots,L_k:A\to\D(O\times R)$, writing
$m(i):=\sum_aq(a)P(i\mid a)$,
\[
 V\bigl(q,P\triangleright(K_1,\ldots,K_k)\bigr)
 -V\bigl(q,P\triangleright(L_1,\ldots,L_k)\bigr)
 =\sum_{i:m(i)>0}m(i)
 \bigl[V(q_i,K_i)-V(q_i,L_i)\bigr].
\]
Hence every $\alpha V+\beta$, with $\alpha>0$, is common-scale and branch-additive.

Conversely, let the branch-additive common-scale representation be $W=\varphi\circ V$ and define
\[
 \psi(x):=\varphi(\underline u+x)-\varphi(\underline u),
 \qquad x\ge0.
\]
Fix $n\ge2$ and let $q$ be uniform on $n$ actions.  For $x\in[0,\lambda\log n]$, choose a constant-$o_*$ erasure
channel whose value is $\underline u+x$.  In the first compound, use this channel after the first branch
of a public coin with probabilities $\theta$ and $1-\theta$, and use the uninformative constant-$o_*$ channel after
the second.  In the comparison compound, use the uninformative channel after both branches.  The two compound values
are $\underline u+\theta x$ and $\underline u$, so \eqref{eq:branch-additive-index} gives
\[
 \psi(\theta x)=\theta\psi(x)
 \qquad(0\le\theta\le1).
\]
Taking $x=\lambda\log n$ shows that $\psi$ is linear on $[0,\lambda\log n]$.  These intervals are nested and
unbounded, so $\psi(x)=\alpha x$ on $[0,\infty)$ for one $\alpha>0$.  Therefore
$\varphi(v)=\alpha v+\beta$ on the range of $V$.

Finally, any second trace-tempered index $\widetilde V$ is common-scale, because expected utility and mutual
information are invariant under block embedding, and is branch-additive by the same two chain rules.  Applying
part \textup{(ii)} and then restricting to deterministic payoff outcomes and to constant-payoff record channels gives
$\widetilde u=\alpha u+\beta$ and $\widetilde\lambda=\alpha\lambda$.
\end{proof}

\begin{proof}[Proof of Corollary~\ref{cor:signduality}]
Part \textup{(i)} is Theorem~\ref{thm:main}.  For part \textup{(ii)}, reverse
the primitive relation in every environment.  The reversed family satisfies
Axioms~\textup{(A1)}, \textup{(A2)}, \textup{(A5)}, and \textup{(A8)}; Axiom~\textup{(A3)} holds with
its two witnessing outcomes interchanged.  It also satisfies the forward
versions of Axioms~\textup{(A4)}, \textup{(A6)}, and
\textup{(A7)}.  Theorem~\ref{thm:main} therefore represents the reversed
family by
\(
  \E[\widehat u(O)]+\widehat\lambda I(A;O,R)
\)
with \(\widehat\lambda>0\).  Negating that index represents the original
family, so it has the asserted form with
\(u=-\widehat u\) and \(\lambda=-\widehat\lambda<0\).

For part \textup{(iii)}, collapse the action alphabet of any pair $(q,K)$ to
a singleton.  Axiom~$(\mathrm{A}7^0)$ makes the pair indifferent to the
resulting one-action channel, whose output law is $p_{qK}$.  Axiom
$(\mathrm{A}6^0)$ then erases its record without changing value.  Thus every
pair is indifferent to the one-action channel that delivers precisely its
payoff marginal.  Axiom~\textup{(A5)} makes these reductions coherent across
blocks.  On payoff lotteries, Axiom~\textup{(A8)} applied to an
action-independent public coin gives mixture independence, Axiom
\textup{(A2)} gives continuity, and Axiom~\textup{(A3)} gives
non-triviality.  The mixture-space theorem therefore yields a nonconstant
expected-utility index.  Conversely, expected utility is unaffected by record
or action processing and satisfies the three indifference clauses.  The
positive and negative formulae satisfy their respective clauses by data
processing and the two chain rules.
\end{proof}

\subsection{Independence of the axioms}

\begin{proof}[Proof of Proposition~\ref{prop:independence}]
Fix a nonconstant $u$ and write $U(q,K):=\E_{qK}[u(O)]$, $Z:=(O,R)$, and $I:=\I(A;Z)$.  All scores below are
extended to block and compound channels by the same formula.

\emph{Failure of Axiom \textup{(A1)}.}
Use the Pareto relation on the two coordinates $(U,I)$:
\[
 (q,K)\succeq(p,L)
 \quad\Longleftrightarrow\quad
 U(q,K)\ge U(p,L)\ \text{and}\ I(q,K)\ge I(p,L).
\]
It is transitive but incomplete whenever one alternative has more payoff and the other more information.  Its graph is
closed.  Both data-processing axioms are coordinatewise monotonic.  Under a common-branch replacement, both the
payoff and information differences are multiplied by the same positive branch probability, so Axiom~\textup{(A8)}
holds in both directions.  Block coherence holds.  For Axiom~\textup{(A3)}, choose the two outcomes so that
$u(o^+)>u(o^-)$; for Axiom~\textup{(A4)}, the two test lotteries have equal payoff value and information values
$\log2$ and zero.  Thus both relevance axioms hold.

\emph{Failure of Axiom \textup{(A2)}.}
Order alternatives lexicographically by $(U,I)$, with payoff first.  This is a weak order and satisfies all structural,
processing, and continuation requirements.  Axiom~\textup{(A3)} holds through the first coordinate, and
Axiom~\textup{(A4)} through the second when the first is tied.  It is not continuous even within one fixed channel.  Let a
three-action channel reveal the action fully and let only action $2$ pay the high outcome.  Put
\[
 p=(1/4,1/2,1/4),\qquad q=(0.49,0.50,0.01),\qquad
 q_n=(0.49-\varepsilon_n,0.50+\varepsilon_n,0.01),
\]
where $\varepsilon_n\downarrow0$.  The lotteries $p$ and $q$ have the same expected utility but
$\Sh(q)<\Sh(p)$, so $p\succ_Kq$.  Every $q_n$ has strictly higher expected utility than $p$, so
$q_n\succ_Kp$.  Thus the graph in Axiom~\textup{(A2)} is not closed.

\emph{Failure of Axiom \textup{(A3)}.}
Use $V:=I$.  All other axioms hold.  In particular, the test in
Axiom~\textup{(A4)} has information values $\log2$ and zero.  But the two pure
lotteries in Axiom~\textup{(A3)} both have zero information, whatever their payoffs.

\emph{Failure of Axiom \textup{(A4)}.}
Use $V:=U$.  All other axioms hold.  The outcomes in Axiom~\textup{(A3)} can be
chosen to have different utility, but the two lotteries in Axiom~\textup{(A4)}
always have the same payoff and hence are indifferent.

\emph{Failure of Axiom \textup{(A5)}.}
Treat outer block labels as syntactic action tags.  For an action alphabet presented with at least three outer
blocks, put $w(a^0)=1$ and $w(a^i)=0$ for $i\ne0$; put $w(a)=0$ on every action alphabet with fewer than three
outer blocks.  For $\varepsilon>0$, use within each environment
\[
 V_\varepsilon(q,K):=U(q,K)+I(q,K)+\varepsilon\sum_aq(a)w(a).
\]
This defines a continuous weak order in every fixed environment.  Every comparison in Axioms~\textup{(A6)},
\textup{(A7)}, and \textup{(A8)} is made in a two-block environment, where $w=0$; those axioms therefore reduce
to their valid $U+I$ comparisons.  The fixed relevance tests can be chosen on untagged action sets and also reduce
to $U+I$.  But irrelevant-block coherence fails.  Take three identical one-action, constant-payoff, silent
channels.  Their first two block-supported lotteries are indifferent in the canonical two-block comparison, while
in the three-block environment the lottery supported on block $0$ is strictly preferred to the lottery supported
on block $1$ by $\varepsilon$.

\emph{Failure of Axiom \textup{(A6)}.}
For $\varepsilon>0$, use
\begin{equation}
 V(q,K):=U(q,K)+I(A;Z)-\varepsilon\Sh(Z\mid A).
 \label{eq:no-record-dp}
\end{equation}
This score is continuous and branch-additive, since both mutual information and conditional entropy obey the chain
rule.  Under action processing $A\to B$, the score difference between the original and processed alternatives is
\[
 \bigl[I(A;Z)-I(B;Z)\bigr]
 +\varepsilon\bigl[\Sh(Z\mid B)-\Sh(Z\mid A)\bigr]
 =(1+\varepsilon)I(A;Z\mid B)\ge0,
\]
so \eqref{eq:no-record-dp} satisfies Axiom~\textup{(A7)}.  Axiom~\textup{(A3)} holds through $U$.  In the
Axiom~\textup{(A4)} test, the first lottery has $I=\log2$ and $\Sh(Z\mid A)=0$, while the second has
$I=0$ and $\Sh(Z\mid A)=\log2$; hence the first is strictly preferred.  Record data
processing fails: with constant payoff and a fair record independent of $A$, the original value is $-\varepsilon\log2$;
deleting that noise gives value zero.

\emph{Failure of Axiom \textup{(A7)}.}
Let $G(q):=1-\sum_aq(a)^2$ be Gini entropy and define its posterior reduction
\[
 J_G(q,K):=G(q)-\int G(\pi)\,d\mu_{qK}(\pi),
 \qquad V:=U+J_G.
\]
Concavity of $G$ gives record data processing, and the potential-difference form gives the exact branch chain rule.
Axiom~\textup{(A3)} holds through $U$.  In the Axiom~\textup{(A4)} test, $J_G$ is $1/2$ under the first
lottery and zero under the second, so trace relevance holds.  Action processing can nevertheless increase $J_G$.  Take
$q=(1/8,1/8,3/4)$, a binary record with
\[
 \Pr(0\mid a_1)=\Pr(0\mid a_2)=0,
 \qquad \Pr(0\mid a_3)=1/4,
\]
and merge $a_1,a_2$ into one processed action.  Direct calculation gives $J_G=9/416$ before the merge and
$J_G=3/104$ afterwards.  Hence Axiom~\textup{(A7)} fails.

\emph{Failure of Axiom \textup{(A8)}.}
Use $V:=U+\min\{I,1\}$, with information measured in bits.  It is continuous, respects both data-processing axioms,
and satisfies block coherence.  Axiom~\textup{(A3)} holds through $U$, while the test in Axiom~\textup{(A4)}
has information values one and zero.  Let $q$ be uniform on four actions.  A first-stage record reveals which of two pairs
contains the action.  Hold the continuation following one record fixed and silent.  Following the other record,
keep payoffs constant and compare a silent continuation with one that reveals the action within its pair.  The latter continuation is strictly
preferred at the branch posterior, but the two compounds have respectively one and three-halves bits before capping
and are therefore indifferent.  The forward implication in Axiom~\textup{(A8)} fails.
\end{proof}

\subsection{Proofs of the choice implications}\label{app:choice-proofs}

Several results stated only in summary form in the text are given here in full.  Proofs appear in an order that
front-loads shared machinery.  Throughout this subsection, write
$Z:=O\times R$ and $K_a:=K(\cdot\mid a)\in\D(Z)$ for the row of action $a$; the distribution induced by $q$ is
$p_{qK}$, as in Section~\ref{sec:theorem}.

\begin{proof}[Proof of Proposition~\ref{prop:signature}]
(i) Under either lottery, each visible pair $(o_0,r_j)$, $j=1,2$, has probability $\tfrac12$: under $q'$ because
each $a_j$ produces $(o_0,r_j)$ for sure, under $q''$ because each of $a_3,a_4$ produces either pair with equal
probability.  Hence $p_{q'K}=p_{q''K}$.  On $\supp(q')$ the rows are degenerate and distinct, so the visible pair
reveals the action and $\I_{q'K}=\Sh(q')=\log2$; on $\supp(q'')$ the rows are identical, so the visible pair is
independent of the action and $\I_{q''K}=0$.  The payoff terms agree because the payoff marginals do, giving the
displayed gap.
(ii) All rows of $K^\circ$ are equal, so $\I_{qK^\circ}=0$ for every $q$; since the payoff is constantly $o_0$,
the representation gives indifference.
\end{proof}

\begin{proof}[Proof of Proposition~\ref{prop:optimal-choice}]
Expected utility is linear in $q$, while mutual information is concave in the input distribution of a fixed channel
\cite{CoverThomas2006}.  Continuity on the compact simplex gives existence.  Using
\[
 I_{qK}(A;Z)=\sum_aq(a)\KL\!\left(K_a\,\middle\|\,p_{qK}\right),
\]
its directional derivative in the mass on action $a$ equals the displayed divergence up to a term common to every
action.  The result is the Kuhn--Tucker condition for a concave programme.  At constant $u$ it reduces to the classical
capacity condition.
\end{proof}

\begin{proposition}[Optimal marginal and multiplicity]\label{prop:optimal-marginal}
Let $\lambda>0$ and write $C^*(K):=\max_{q\in\D(A)}V(q,K)$.  There is a unique $p^*\in\D(Z)$ such that
$p_{q^*K}=p^*$ for every optimal lottery $q^*$.  With
\[
 \mathcal A^*
 :=\bigl\{a\in A:
 \bar u(a)+\lambda
 \KL\!\left(K_a\,\middle\|\,p^*\right)=C^*(K)
 \bigr\},
\]
the set of optimal lotteries is exactly
\[
 \bigl\{q\in\D(A):
 p_{qK}=p^*,\ \supp(q)\subseteq\mathcal A^*
 \bigr\}.
\]
The optimum is unique whenever the rows $\{K_a:a\in\mathcal A^*\}$ are affinely independent.  If actions are
partitioned into classes with identical rows, $V(q,K)$ depends on $q$ within each class only through the class's
total mass.  The optimal lotteries are exactly the lifts of the optimal lotteries of the collapsed channel, and
each class's mass may be distributed arbitrarily among its members.
\end{proposition}

\begin{proof}[Proof of Proposition~\ref{prop:optimal-marginal}]
The entropy identity gives
\[
 V(q,K)
 =\lambda\Sh(p_{qK})
  +\sum_aq(a)\bigl[\bar u(a)-\lambda\Sh(K_a)\bigr].
\]
If $q_0$ and $q_1$ are optimal with $p_{q_0K}\ne p_{q_1K}$, strict concavity of Shannon entropy and linearity of
$q\mapsto p_{qK}$ and of the remaining sum give $V(tq_0+(1-t)q_1,K)>C^*(K)$ for $t\in(0,1)$, a contradiction.
Hence every optimum induces the same $p^*$.  By Proposition~\ref{prop:optimal-choice}, an optimum $q^*$ satisfies
the displayed equalities at some constant $c$; averaging them under $q^*$ gives $c=V(q^*,K)=C^*(K)$, so
$\supp(q^*)\subseteq\mathcal A^*$.  Conversely, if $p_{qK}=p^*$ and $\supp(q)\subseteq\mathcal A^*$, then
\[
 V(q,K)
 =\sum_aq(a)\bigl[\bar u(a)+\lambda\KL(K_a\|p^*)\bigr]
 =C^*(K).
\]
Affine independence makes the representation $p^*=\sum_aq(a)K_a$ with $\supp(q)\subseteq\mathcal A^*$ unique.
Identical rows share $\bar u(a)$ and $\Sh(K_a)$, so both $p_{qK}$ and the displayed expression for $V$ are
unchanged by redistribution within a class; collapsing the classes and lifting aggregate masses proves the last
claim.
\end{proof}

\begin{proposition}[Support and pure choice]\label{cor:support-pure}
Let $\lambda>0$ and let $Z_K:=\{z\in Z:K_a(z)>0\text{ for some }a\in A\}$.
\begin{enumerate}[label=\textup{(\roman*)},leftmargin=*]
\item Every optimal lottery $q^*$ satisfies $p_{q^*K}(z)>0$ for every $z\in Z_K$.  Consequently, if every action
$a$ has a private consequence---some $z_a$ with $K_a(z_a)>0$ and $K_b(z_a)=0$ for all $b\ne a$---then every
optimal lottery has full support.
\item The pure lottery $\delta_a$ is optimal if and only if
\[
 \bar u(b)+\lambda\KL\!\left(K_b\,\middle\|\,K_a\right)
 \le \bar u(a)
 \qquad\text{for every }b\in A,
\]
and uniquely optimal if these inequalities are strict for $b\ne a$.  In particular, if some rival reaches a
consequence outside the support of $K_a$, then $\delta_a$ is not optimal, however large its payoff advantage.
\end{enumerate}
\end{proposition}

\begin{proof}[Proof of Proposition~\ref{cor:support-pure}]
Throughout, $p^*$ denotes the common optimal marginal of Proposition~\ref{prop:optimal-marginal}.
(i) Suppose $p^*(z)=0$ for some $z\in Z_K$ and choose $a$ with $K_a(z)>0$.  Then
$\KL(K_a\|p^*)=+\infty$, which is incompatible with the finite on-support equality if some optimum uses $a$ and
with the off-support inequality of Proposition~\ref{prop:optimal-choice} if none does.  Under the
private-consequence condition, an optimum with $q^*(a)=0$ would force $p^*(z_a)=0$.

(ii) For $q=\delta_a$ the induced marginal is $K_a$ and the on-support constant is $\bar u(a)$; the displayed
system is the off-support condition of Proposition~\ref{prop:optimal-choice}, necessary and sufficient by
concavity.  Under strict inequalities, $\mathcal A^*=\{a\}$ in Proposition~\ref{prop:optimal-marginal}, so every
optimum is $\delta_a$.
\end{proof}

\begin{proposition}[Optimal mixing]\label{prop:binary}
Let $\lambda>0$ and $A=\{a,b\}$ with $K_a\ne K_b$ and
$\Delta:=\bar u(a)-\bar u(b)\ge0$.  For $x\in[0,1]$ let $q_x$ put mass $x$ on $b$, so that
$p_{q_xK}=(1-x)K_a+xK_b=:p_x$, and for $x\in(0,1)$ define
\[
 g(x):=\KL\!\left(K_b\,\middle\|\,p_x\right)-\KL\!\left(K_a\,\middle\|\,p_x\right).
\]
\begin{enumerate}[label=\textup{(\roman*)},leftmargin=*]
\item $g$ is continuous and strictly decreasing, with $g(0^+)=\KL(K_b\,\|\,K_a)$ and
$g(1^-)=-\KL(K_a\,\|\,K_b)$ as extended-real limits.  The pure lottery $\delta_a$ is optimal if and only if
$\Delta\ge\lambda\,\KL(K_b\,\|\,K_a)$.  Below this threshold the optimum is unique and interior,
$q^*=q_{x^*}$ with $x^*$ the unique solution of $\Delta=\lambda g(x)$.
\item The optimal mass $x^*(\Delta,\lambda)$ on the inferior action is nonincreasing in $\Delta$ and
nondecreasing in $\lambda$, strictly so on the mixing region; at $\Delta=0$ it is the capacity-achieving input
of the two-row channel, independently of $\lambda$.
\end{enumerate}
\end{proposition}

\begin{proof}
Write $f(x):=V(q_x,K)=\bar u(a)-x\Delta+\lambda I(x)$ with
$I(x)=\Sh(p_x)-(1-x)\Sh(K_a)-x\Sh(K_b)$.  Since $K_a\ne K_b$,
$x\mapsto p_x$ is affine and injective, so strict concavity of $\Sh$ \cite{CoverThomas2006} makes $f$ strictly
concave on $[0,1]$, and a unique maximiser exists.  On $(0,1)$ every $z\in\supp K_a\cup\supp K_b$ has $p_x(z)>0$
and, using $\sum_z(K_b(z)-K_a(z))=0$,
\[
 I'(x)=\Sh(K_a)-\Sh(K_b)-\sum_z\bigl(K_b(z)-K_a(z)\bigr)\log p_x(z)=g(x),
\]
so $f'(x)=-\Delta+\lambda g(x)$, with $g$ continuous and strictly decreasing as the derivative of a strictly
concave function.  As $x\downarrow0$, $p_x(z)\to K_a(z)$; if $\supp K_b\subseteq\supp K_a$ every summand converges
and $g(0^+)=\KL(K_b\|K_a)$, while if some $\bar z$ has $K_b(\bar z)>0=K_a(\bar z)$ the summand
$-K_b(\bar z)\log\bigl(xK_b(\bar z)\bigr)$ diverges to $+\infty$, matching the extended-sense divergence.
Exchanging the roles of $a$ and $b$ gives the limit at $1$.  By concavity and the mean value theorem,
$\delta_a$ is optimal if and only if $\lim_{x\downarrow0}f'(x)\le0$, i.e.\ $\Delta\ge\lambda\,\KL(K_b\|K_a)$;
this is the threshold stated in the text, and Proposition~\ref{cor:support-pure}\textup{(ii)} with two actions.
If $\Delta<\lambda\,\KL(K_b\|K_a)$, the maximiser is not $x=0$ by the threshold and not $x=1$
because $f'(1^-)=-\Delta-\lambda\KL(K_a\|K_b)<0$; it is therefore the unique root of $\Delta=\lambda g(x)$, which
exists because $g$ decreases strictly between its limits.  This proves \textup{(i)}.  The comparative statics of
\textup{(ii)} follow by inverting the
strictly decreasing $g$ at $\Delta/\lambda$: raising $\Delta$ raises the target, raising $\lambda$ lowers it; at
$\Delta=0$ the condition $g(x)=0$ is the capacity condition of Proposition~\ref{prop:optimal-choice} at constant
payoff.
\end{proof}

\begin{remark}[Overlapping traces and IIA]\label{rem:overlap-iia-example}
Let three actions have equal expected utility and rows
\[
 K_1=(1,0,0),\qquad K_2=(0,1,0),\qquad K_3=(1/2,0,1/2).
\]
With only the first two actions, the optimal lottery is $(1/2,1/2)$.  Adding the third gives
\[
 q^*=(1/3,4/9,2/9),\qquad p_{q^*K}=(4/9,4/9,1/9).
\]
Each active row has divergence $\log(9/4)$ from $p_{q^*K}$, so
Proposition~\ref{prop:optimal-choice} verifies optimality.  The ratio of the first two choice probabilities falls
from $1$ to $3/4$.  Thus overlapping rows can generate the red-bus/blue-bus pattern even when expected utilities
are equal.
\end{remark}

\begin{proposition}[Trace demand]\label{prop:trace-demand}
Fix $K$ and write $e(q):=\E_{qK}[u(O)]$, $i(q):=\I_{qK}(A;O,R)$,
\[
 V^*(\lambda):=\max_{q\in\D(A)}\{e(q)+\lambda i(q)\},
 \qquad
 Q(\lambda):=\arg\max_{q\in\D(A)}\{e(q)+\lambda i(q)\}.
\]
\begin{enumerate}[label=\textup{(\roman*)},leftmargin=*]
\item $V^*$ is convex on $(0,\infty)$, $i$ is affine on the convex compact set $Q(\lambda)$, and
\[
 (V^*)'_-(\lambda)=\min_{q\in Q(\lambda)}i(q),
 \qquad
 (V^*)'_+(\lambda)=\max_{q\in Q(\lambda)}i(q);
\]
at every point of differentiability, $(V^*)'(\lambda)=i(q)$ for every $q\in Q(\lambda)$.
\item Let $0<\lambda_1<\lambda_2$ and $q_j\in Q(\lambda_j)$, writing $e_j:=e(q_j)$, $i_j:=i(q_j)$.  Then
$i_1\le i_2$, $e_1\ge e_2$, and
\[
 \lambda_1(i_2-i_1)\le e_1-e_2\le\lambda_2(i_2-i_1).
\]
Hence either $(e_1,i_1)=(e_2,i_2)$, or $i_1<i_2$ and $e_1>e_2$ and
\[
 \lambda_1\le\frac{e_1-e_2}{i_2-i_1}\le\lambda_2.
\]
\end{enumerate}
\end{proposition}

\begin{proof}[Proof of Proposition~\ref{prop:trace-demand}]
(i) $V^*$ is a pointwise maximum of functions affine in $\lambda$, hence convex and, since
$0\le i\le\log|A|$, finite.  On the convex compact set $Q(\lambda)$ the objective is constant and $e$ is linear, so
$i=(V^*-e)/\lambda$ is affine there.  For $q\in Q(\lambda)$ and $\mu>0$,
$V^*(\mu)\ge V^*(\lambda)+(\mu-\lambda)i(q)$, so each such $i(q)$ is a subgradient.  For the right derivative,
choose $q_h\in Q(\lambda+h)$; then
\[
 \max_{q\in Q(\lambda)}i(q)
 \le\frac{V^*(\lambda+h)-V^*(\lambda)}{h}
 \le i(q_h),
\]
the first inequality by the subgradient bound, the second by optimality of $Q(\lambda)$ at $\lambda$.  By
compactness and continuity, accumulation points of $(q_h)_{h\downarrow0}$ lie in $Q(\lambda)$, so the right-hand
side accumulates in $i(Q(\lambda))$ and the right derivative equals the maximum.  The left derivative is symmetric,
and a convex function on an interval is differentiable off a countable set.

(ii) Optimality at the two weights gives
$e_1+\lambda_1i_1\ge e_2+\lambda_1i_2$ and $e_2+\lambda_2i_2\ge e_1+\lambda_2i_1$.  Adding yields
$(\lambda_2-\lambda_1)(i_2-i_1)\ge0$, hence $i_2\ge i_1$; the two inequalities rearrange to the displayed bracket,
which gives $e_1\ge e_2$ and, when the pairs differ, the final bound.
\end{proof}

\begin{corollary}[Extreme trace weights]\label{cor:lambda-limits}
Let $e^0:=\max_qe(q)$, $i^0:=\max\{i(q):e(q)=e^0\}$, $C:=\max_qi(q)$, and $e^C:=\max\{e(q):i(q)=C\}$.  If
$\lambda_n\downarrow0$ and $q_n\in Q(\lambda_n)$, every accumulation point of $(q_n)$ maximises $i$ among the
maximisers of $e$; if $\lambda_n\uparrow\infty$, every accumulation point maximises $e$ among the maximisers of
$i$.  Moreover
\[
 V^*(\lambda)=e^0+\lambda i^0+o(\lambda)
 \quad\text{as }\lambda\downarrow0,
 \qquad
 V^*(\lambda)=\lambda C+e^C+o(1)
 \quad\text{as }\lambda\uparrow\infty.
\]
\end{corollary}

\begin{proof}[Proof of Corollary~\ref{cor:lambda-limits}]
For $\lambda_n\downarrow0$, comparing $q_n$ with an $e$-maximiser attaining $i^0$ gives
\[
 0\le e^0-e(q_n)\le\lambda_n\bigl(i(q_n)-i^0\bigr).
\]
Boundedness of $i$ forces $e(q_n)\to e^0$, and the same inequality gives $i(q_n)\ge i^0$; continuity delivers the
accumulation claim.  Dividing by $\lambda_n$,
\[
 \frac{V^*(\lambda_n)-e^0}{\lambda_n}=i(q_n)-\frac{e^0-e(q_n)}{\lambda_n},
\]
where the last fraction lies in $[0,\,i(q_n)-i^0]$ and vanishes, proving the first expansion.  For
$\lambda_n\uparrow\infty$, comparing $q_n$ with an $i$-maximiser attaining $e^C$ gives
$0\le\lambda_n(C-i(q_n))\le e(q_n)-e^C$, and the symmetric argument completes the proof.
\end{proof}

\begin{proof}[Proof of Proposition~\ref{prop:garbling-price}]
The processor leaves the payoff coordinate unchanged, so the expected-utility terms coincide and the difference is
$\lambda[\I_{qK}(A;O,R)-\I_{qL}(A;O,R')]$.  Write $Z:=(O,R)$ and $Z':=(O,R')$.  By construction $A\to Z\to Z'$ is
a Markov chain, so expanding $\I(A;Z,Z')$ in both orders gives
\[
 \I(A;Z)-\I(A;Z')=\I(A;Z\mid Z')=\I(A;R\mid O,R'),
\]
the last equality because $O$ is part of $Z'$.  Conditional mutual information is nonnegative and vanishes exactly
when $A$ and $(O,R)$ are conditionally independent given $(O,R')$ \cite{CoverThomas2006}.  This is the zero-price
condition stated in the proposition.  Under a garbling, it is also equivalent to equality of the posterior laws,
$\mu_{qK}=\mu_{qL}$.  Indeed, the Markov property gives
$\Pr(A=\cdot\mid Z,Z')=\pi^{qK}_{Z}$ almost surely, while conditional independence gives
$\Pr(A=\cdot\mid Z,Z')=\pi^{qL}_{Z'}$; the two random posteriors then coincide almost surely.  Conversely, by the
identity in the proof of Corollary~\ref{cor:matching}, equality of the posterior laws gives equality of the
information terms, and the payoff terms already agree.
\end{proof}

\begin{corollary}[The recorded coin]\label{cor:recorded-coin}
Let $K,L:A\to\D(O\times R)$ and $t\in(0,1)$.  Let $tK\oplus(1-t)L$ be the public mixture that draws an
action-independent coin $Y$ with $\Pr(Y=0)=t$, applies $K$ if $Y=0$ and $L$ otherwise, and includes the coin in
the record; let $M_t:=tK+(1-t)L$ be the rowwise mixture, which hides it.  Then, for every $q\in\D(A)$,
\begin{enumerate}[label=\textup{(\roman*)},leftmargin=*]
\item $V\bigl(q,\,tK\oplus(1-t)L\bigr)=t\,V(q,K)+(1-t)\,V(q,L)$;
\item $V\bigl(q,\,tK\oplus(1-t)L\bigr)-V(q,M_t)=\lambda\,\I(A;Y\mid O,R)\ge0$, with equality if and only if $A$
and $Y$ are conditionally independent given $(O,R)$.
\end{enumerate}
\end{corollary}

\begin{proof}[Proof of Corollary~\ref{cor:recorded-coin}]
Part \textup{(i)} is the chain rule with an action-independent first stage:
$\I(A;Y)=0$, and each reached branch reproduces the joint law of the
corresponding constituent.  Part \textup{(ii)} follows from
Proposition~\ref{prop:garbling-price}: deleting the coin coordinate is a
deterministic payoff-preserving record processor carrying the public mixture
to $M_t$.
\end{proof}

\begin{definition}[Garbling and more-capable orders]\label{def:channel-orders}
Let $K:A\to\D(O\times R)$ and $L:A\to\D(O\times R')$ have the same payoff kernels.  Write $K\succeq_gL$ if
$L=KT$ for some record processor $T$, and write $K\succeq_cL$ if
$\I_{qK}(A;O,R)\ge\I_{qL}(A;O,R')$ for every $q\in\D(A)$; the latter is the \emph{more-capable} order on the
visible-consequence channels \cite{KornerMarton1977}.
\end{definition}

\begin{proposition}[Uniform dominance is more-capable, not garbling]\label{prop:blackwell}
Let $K$ and $L$ have identical payoff kernels.
\begin{enumerate}[label=\textup{(\roman*)},leftmargin=*]
\item $K\succeq_cL$ if and only if $V(q,K)\ge V(q,L)$ for every $q\in\D(A)$.
\item $K\succeq_gL\Rightarrow K\succeq_cL$, but the converse fails.  There are channels for which
$K\succeq_cL$ although $L\ne KT$ for every record processor $T$.
\end{enumerate}
\end{proposition}

\begin{proof}[Proof of Proposition~\ref{prop:blackwell}]
(i) Because the payoff kernels coincide,
$V(q,K)-V(q,L)=\lambda[\I_{qK}-\I_{qL}]$ with $\lambda>0$.  Uniform value dominance and
$K\succeq_cL$ are therefore the same statement.

(ii) If $L=KT$, Proposition~\ref{prop:garbling-price} gives $\I_{qK}\ge\I_{qL}$ at every $q$, and hence
$K\succeq_gL\Rightarrow K\succeq_cL$.  For failure of the converse, let $A=\{0,1\}$, fix $o^\ast\in O$, and let
both channels pay $o^\ast$ surely.  Let $K$ carry the binary
erasure record with erasure probability $\varepsilon=\tfrac12$ and $L$ the binary symmetric record with crossover
$p=\tfrac15$:
\[
 K(o^\ast,a\mid a)=1-\varepsilon,\quad K(o^\ast,\mathrm e\mid a)=\varepsilon,
 \qquad
 L(o^\ast,a\mid a)=1-p,\quad L(o^\ast,1-a\mid a)=p.
\]
The payoff kernels agree row by row, so it suffices to compare informations, and the comparison is invariant to the
base; compute in bits with $h$ the binary entropy.  For $q=(1-t,t)$, grouping records gives
\[
 \I_{qK}=(1-\varepsilon)\,h(t),
 \qquad
 \I_{qL}=h(p\ast t)-h(p),
 \qquad p\ast t:=p(1-t)+(1-p)t.
\]
By Mrs Gerber's Lemma \cite{WynerZiv1973}, $v\mapsto h\bigl(p\ast h^{-1}(v)\bigr)$ is convex on $[0,1]$, with
$h^{-1}$ valued in $[0,\tfrac12]$; since $h(p\ast t)$ is invariant under $t\mapsto1-t$, take $t\le\tfrac12$ and
$v=h(t)$.  The chord between $v=0$ and $v=1$ gives
\[
 h(p\ast t)\le(1-v)\,h(p)+v,
 \qquad\text{hence}\qquad
 \I_{qL}\le\bigl(1-h(p)\bigr)h(t)\le(1-\varepsilon)\,h(t)=\I_{qK},
\]
the last step because $\varepsilon=\tfrac12\le h(\tfrac15)$.  Thus $K\succeq_cL$.  Now suppose $L=KT$ for a
processor $T$, and write $t_r:=T(0\mid o^\ast,r)$ for $r\in\{0,1,\mathrm e\}$.  Matching the probability of the
processed record $0$ in the two rows,
\[
 (1-\varepsilon)t_0+\varepsilon t_{\mathrm e}=1-p,
 \qquad
 (1-\varepsilon)t_1+\varepsilon t_{\mathrm e}=p,
\]
and subtracting gives $t_0-t_1=(1-2p)/(1-\varepsilon)=\tfrac65>1$, impossible.  Compositions of record processors
are record processors, so no chain of garblings produces $L$ from $K$ either.
\end{proof}

The next definition and proposition formalise the side-bet elicitation of \S\ref{sec:measurement}.  Let $o^+$
and $o^-$ attain the maximum and minimum of $u$, write $\Delta u:=u(o^+)-u(o^-)>0$, and for $\beta\in[0,1]$ let
$w(\beta):=\beta u(o^+)+(1-\beta)u(o^-)$.

\begin{definition}[Side bet]\label{def:dilution}
For $K:A\to\D(O\times R)$, put $R^\ast:=R\sqcup\{\ast\}$ and regard $K$ as a channel on $R^\ast$ by zero
padding.  For $\varepsilon\in[0,1)$ and $\beta\in[0,1]$, let $P_\varepsilon:A\to\D(\{1,2\})$ be the public coin
$P_\varepsilon(2\mid a)=\varepsilon$, and let $B_\beta:A\to\D(O\times R^\ast)$ have the constant rows
$B_\beta(o^+,\ast\mid a)=\beta$ and $B_\beta(o^-,\ast\mid a)=1-\beta$.  Write
\[
 D_{\varepsilon,\beta}K:=P_\varepsilon\triangleright(K,B_\beta),
\]
as in \eqref{eq:compound}.
\end{definition}

\begin{proposition}[Willingness to pay for trace]\label{prop:wtp}
Fix $q\in\D(A)$, $p\in\D(B)$, and channels $K:A\to\D(O\times R)$ and $L:B\to\D(O\times S)$.
\begin{enumerate}[label=\textup{(\roman*)},leftmargin=*]
\item For every $\varepsilon\in[0,1)$ and $\beta\in[0,1]$,
\[
 V\bigl(q,\,D_{\varepsilon,\beta}K\bigr)=(1-\varepsilon)\,V(q,K)+\varepsilon\,w(\beta).
\]
\item If $\varepsilon\in(0,1)$ and $\beta,\beta'\in[0,1]$ satisfy
$\bigl(q,D_{\varepsilon,\beta}K\bigr)\sim\bigl(p,D_{\varepsilon,\beta'}L\bigr)$, then
\[
 \frac{\varepsilon}{1-\varepsilon}\bigl[w(\beta)-w(\beta')\bigr]
 =V(p,L)-V(q,K).
\]
When the material parts agree, $\E_{qK}[u(O)]=\E_{p,L}[u(O)]$, the left-hand side equals
$\lambda\bigl[\I_{pL}(B;O,S)-\I_{qK}(A;O,R)\bigr]$: information is priced linearly at rate $\lambda$.
\item Such a pair $(\beta,\beta')$ exists if and only if
$\lvert V(p,L)-V(q,K)\rvert\le\frac{\varepsilon}{1-\varepsilon}\Delta u$; a boundary response therefore brackets
the value difference instead of measuring it.
\end{enumerate}
\end{proposition}

\begin{proof}
(i) The material part of the compound is $(1-\varepsilon)\E_{qK}[u(O)]+\varepsilon w(\beta)$.  For the trace part,
the coin is action-independent, so both branch posteriors equal $q$ and the chain rule for the compound gives an
information term of $(1-\varepsilon)\I_{qK}(A;O,R)$, the side-bet branch contributing zero.
(ii) By the moreover clause of Theorem~\ref{thm:main}, the indifference holds if and only if
$(1-\varepsilon)V(q,K)+\varepsilon w(\beta)=(1-\varepsilon)V(p,L)+\varepsilon w(\beta')$; rearrange, and
substitute the representation when the material parts cancel.
(iii) $w(\beta)-w(\beta')$ is continuous on $[0,1]^2$ with range $[-\Delta u,\Delta u]$; apply the intermediate
value theorem to the equation in \textup{(ii)}.
\end{proof}

\begin{lemma}[The information--payoff plane]\label{lem:info-payoff-plane}
Write $\phi(q,K):=\bigl(\I_{qK}(A;O,R),\,E_{qK}\bigr)$, and let $\underline u:=\min_ou(o)$ and
$\overline u:=\max_ou(o)$.  The image of $\phi$ over all finite pairs $(q,K)$ is exactly
$[0,\infty)\times[\underline u,\overline u]$, and the indifference classes of the block-comparison relation are,
in these coordinates, the nonempty intersections of the lines $E=c-\lambda\I$ with that strip.
\end{lemma}

\begin{proof}
Inclusion is immediate.  Conversely, fix $(t,w)$ in the strip, choose $n$ with $\log n\ge t$, and on
$A=\{1,\dots,n\}$ use continuity of entropy along a segment from a point mass to the uniform lottery to find
$q_t$ with $\Sh(q_t)=t$.  Let $K(o^+,a\mid a)=\beta$ and $K(o^-,a\mid a)=1-\beta$ with $\beta$ chosen so that
$w(\beta)=w$.  The payoff lottery is action-independent with mean $w$, and the record reveals the action, so
$\phi(q_t,K)=(t,w)$.  The description of the indifference classes then reads the representation through $\phi$,
using the moreover clause of Theorem~\ref{thm:main}.
\end{proof}

\begin{proof}[Proof of Proposition~\ref{prop:single-price}]
The identity rearranges $E_{qK}+\lambda\I_{qK}=E_{pL}+\lambda\I_{pL}$.  For uniqueness, let
$\lambda'\ne\lambda$, put $t:=\tfrac12\min\{1,\,(\overline u-\underline u)/\lambda\}>0$, and use
Lemma~\ref{lem:info-payoff-plane} to realise pairs with $\phi$-coordinates $(t,\overline u-\lambda t)$ and
$(0,\overline u)$; they are indifferent, while $E+\lambda'\I$ differs across them by $(\lambda-\lambda')t\ne0$.
\end{proof}

To compare agents, let both families be defined over the same outcome set $O$, on the common domain of all finite
pairs, with mutual information in one fixed base.  A
channel $B$ is \emph{silent} if all its rows are equal; then the joint law under any $p$ is a product, so
$\I_{pB}=0$ and $V^i(p,B)$ is the expectation of $u_i$ under the $O$-marginal of the common row.  Since
one-action channels realise every payoff lottery in $\D(O)$, silent pairs realise exactly the payoff lotteries.

\begin{definition}[More trace-seeking]
Family 2 is \emph{more trace-seeking than} family 1 if
\[
 (q,K)\succeq^1(p,B)\quad\Longrightarrow\quad(q,K)\succeq^2(p,B)
\]
for all lotteries $q,p$, all channels $K$, and all silent channels $B$.
\end{definition}

\begin{proposition}[More trace-seeking]\label{prop:comparative}
Let $\{\succeq^1_K\}$ and $\{\succeq^2_K\}$ satisfy Axioms~\textup{(A1)--(A8)}.  Then family 2 is more
trace-seeking than family 1 if and only if the representations can be chosen with a common utility and
$\lambda_2\ge\lambda_1$.
\end{proposition}

\begin{proof}[Proof of Proposition~\ref{prop:comparative}]
Throughout, block comparisons are read through the moreover clause of Theorem~\ref{thm:main}.

Suppose the representations are $(u,\lambda_1)$ and $(u,\lambda_2)$ with $\lambda_2\ge\lambda_1$.  For any
$(q,K)$ and silent $(p,B)$,
\[
 V^2(q,K)-V^2(p,B)
 =\bigl[V^1(q,K)-V^1(p,B)\bigr]+(\lambda_2-\lambda_1)\I_{qK}
 \ge V^1(q,K)-V^1(p,B),
\]
since $\I_{qK}\ge0$ and the silent pair's value is material under both.  A weak (or strict) agent-1 preference
therefore transfers to agent 2, and reading the display from right to left gives the mirrored clause
$(p,B)\succeq^2(q,K)\Rightarrow(p,B)\succeq^1(q,K)$: the weak one-sided definition, its strict variant, and the
two-sided variant coincide on this class.

Conversely, suppose $\{\succeq^2_K\}$ is more trace-seeking than $\{\succeq^1_K\}$, with representations
$(u_1,\lambda_1)$ and $(u_2,\lambda_2)$.

\emph{Step 1: common utility.}  Let $D:=\{\mu-\nu:\mu,\nu\in\D(O)\}$ and $f_i(d):=\sum_ou_i(o)d(o)$; each
$f_i\ne0$ on $D$ because $u_i$ is nonconstant.  Silent pairs realise every payoff lottery, and the definition's
left argument ranges over all pairs, so silent-against-silent instances are in its scope:
$f_1(d)\ge0\Rightarrow f_2(d)\ge0$ on $D$.  If $f_1(d)=0$, applying this to $d$ and $-d$ gives $f_2(d)=0$, so
$\ker f_1\subseteq\ker f_2$ and $f_2=\alpha f_1$ for some $\alpha\in\mathbb R$.  If $\alpha=0$ then $u_2$ is
constant, contradicting Axiom~\textup{(A3)} for family 2.  (The case is not idle: an index with constant $u_2$,
which violates only that axiom, is indifferent across all silent pairs and is more trace-seeking than every
family whatsoever.)  If $\alpha<0$, pick $o^+,o^-$ with $u_1(o^+)>u_1(o^-)$: the silent instance
$\delta_{o^+}\succ^1\delta_{o^-}$ forces $\delta_{o^+}\succeq^2\delta_{o^-}$, while $\alpha<0$ gives the strict
reverse.  Hence $\alpha>0$ and $u_2=\alpha u_1+\beta$.  Since $((u_2-\beta)/\alpha,\lambda_2/\alpha)$ again
represents family 2, choose $u_2=u_1=:u$, replacing $\lambda_2$ by $\lambda_2/\alpha>0$.  The choice is
innocuous: any two common-utility normalisations rescale both weights by one positive factor, because
$u$ nonconstant forces the factors to agree, so the order of the weights does not depend on it.

\emph{Step 2: weight ordering.}  Write $\overline u:=\max_ou(o)$ and $\underline u:=\min_ou(o)$;
$\overline u>\underline u$ by Axiom~\textup{(A3)}.  Suppose $\lambda_2<\lambda_1$ and fix
$\lambda'\in(\lambda_2,\lambda_1)$.  Choose $t\in\bigl(0,(\overline u-\underline u)/\lambda'\bigr)$ and
$w:=\underline u+\lambda't\in(\underline u,\overline u)$.  By Lemma~\ref{lem:info-payoff-plane} there is a pair
$(q,K)$ with $\I_{qK}=t$ and $E_{qK}=\underline u$; its payoff lottery is action-independent, so after Step~1
both families assign it the same coordinates.  Choose a silent $(p,B)$ whose common-row $O$-marginal pays
$o^+$ with probability $(w-\underline u)/(\overline u-\underline u)$ and $o^-$ otherwise, so
$V^1(p,B)=V^2(p,B)=w$.  Then $V^1(q,K)=\underline u+\lambda_1t>w$, so $(q,K)\succ^1(p,B)$, and the definition
gives $(q,K)\succeq^2(p,B)$; but $V^2(q,K)=\underline u+\lambda_2t<w$, a contradiction.  Hence
$\lambda_2\ge\lambda_1$.
\end{proof}

\begin{remark}
The proof uses $\lambda_i>0$ only in the choice $E_{qK}=\underline u$; taking $E_{qK}=\overline u$ when
$\lambda'<0$ runs both directions verbatim.  The equivalence therefore holds across all three orientations of
Corollary~\ref{cor:signduality}: more trace-seeking and less trace-averse are one comparative.
\end{remark}

\begin{remark}[Eliciting $\lambda$]\label{rem:elicitation}
Normalise $u(o^+)=1$ and $u(o^-)=0$.  All comparisons below are block comparisons, hence within the primitive
domain, and no information quantity is estimated: the information gap is an accounting identity of the design.
\emph{Step 1.}  For each $o\in O$, elicit $\beta_o$ with $(\delta_\ast,K_o)\sim(\delta_\ast,B_{\beta_o})$, where
$K_o$ is the degenerate one-action channel used in Lemma~\ref{lem:relevance-bridge}; one-action channels carry zero information, so
$u(o)=\beta_o$.
\emph{Step 2.}  Fix $o\in O$, $n\ge2$, the uniform $q$ on $\{1,\dots,n\}$, a side-bet probability $\varepsilon\in(0,1)$, and
the constant-payoff channels $K^{\mathrm{id}}_o$ and $K^0_o$ used in Lemma~\ref{lem:relevance-bridge}.  Elicit $\beta$ with
\[
 \bigl(q,\,D_{\varepsilon,\beta}K^{0}_o\bigr)\sim\bigl(q,\,D_{\varepsilon,0}K^{\mathrm{id}}_o\bigr);
 \qquad\text{then}\qquad
 \lambda=\frac{\varepsilon}{1-\varepsilon}\cdot\frac{\beta}{\log n},
\]
by Proposition~\ref{prop:wtp}\textup{(i)} applied to $V(q,K^{\mathrm{id}}_o)=u(o)+\lambda\log n$ and
$V(q,K^{0}_o)=u(o)$.  Such a $\beta$ exists if and only if $\lambda\log n\le\varepsilon/(1-\varepsilon)$; a
subject still strictly preferring the right block at $\beta=1$ reveals
$\lambda>\varepsilon/\bigl((1-\varepsilon)\log n\bigr)$, and $\varepsilon$ is raised.
\emph{Controls.}  Orientation: Theorem~\ref{thm:main} implies
$(q,K^{\mathrm{id}}_o)\succ(q,K^{0}_o)$ at every payoff, although Axiom~\textup{(A4)} assumes the comparison at only one.  Garbling:
under the total record garbling of $K^{\mathrm{id}}_o$, both blocks carry zero information and indifference must
obtain; a surviving strict preference indicates that an uncontrolled feature of the display still traces the
action.  Transport: repeating Step 2 at another $n$, prior, or payoff level must return the same $\lambda$, by
Proposition~\ref{prop:single-price}; an estimate that declines in $\log n$ is the signature of the
capped-information criterion.
\end{remark}
